% !TEX root = ../main.tex
% ---------------------------------------------------------------
% GENERATED FILE - DO NOT EDIT.
% Produced by docs/tools/render_reference.py from the repository source.
% Regenerate with: make docs
% ---------------------------------------------------------------

\apibreak
\section{\texttt{IssueClusterOutput}}
\label{class:feedback_intelligence_agent.tools.IssueClusterOutput}\label{schema:IssueClusterOutput}\index{Classes!IssueClusterOutput}

Recurring issue clusters detected in the feedback dataset.
\apifield{Inherits}\texttt{ToolOutput}
\apifield{Fields}
\begin{arglist}
  \item[\texttt{total\_\allowbreak{}records}] \texttt{int}, required
  \item[\texttt{clusters}] \texttt{list[IssueCluster]\allowbreak{}}, required
\end{arglist}
\apifield{Source}\texttt{src/\allowbreak{}feedback\_\allowbreak{}intelligence\_\allowbreak{}agent/\allowbreak{}tools.\allowbreak{}py:236}~(\href{https://github.com/DiogoRibeiro7/feedback-intelligence-agent/blob/b6cbc710ffdfa5f26f700c999d594c67588c1e16/src/feedback_intelligence_agent/tools.py#L236}{online}), pydantic model, module \cref{module:feedback_intelligence_agent.tools}

\subsection{\texttt{render}}
\label{method:feedback_intelligence_agent.tools.IssueClusterOutput.render}\index{Methods!render}
Summarise the top clusters in one sentence.
\apifield{Usage}
\begin{lstlisting}[style=usage, language=Python]
render(self) -> str
\end{lstlisting}
\apifield{Value}\texttt{str}
\apifield{Source}\texttt{src/\allowbreak{}feedback\_\allowbreak{}intelligence\_\allowbreak{}agent/\allowbreak{}tools.\allowbreak{}py:242}~(\href{https://github.com/DiogoRibeiro7/feedback-intelligence-agent/blob/b6cbc710ffdfa5f26f700c999d594c67588c1e16/src/feedback_intelligence_agent/tools.py#L242}{online})
